% --- Archon physics formalization source begin ---
% archon:physics
% archon:covers PhyXMiniProblems/problem_phyx_mini_0345.lean
% archon:source-report reports/phyx_mini/problem_phyx_mini_0345.source.json

\chapter{Physics problem phyx\_mini\_0345}
\label{ch:PhyXMiniProblems_problem_phyx_mini_0345}

\paragraph{Problem source.}
\#\# Physical scenario

A large research balloon containing \$2.00 \textbackslash{}times 10\textasciicircum{}3 \textbackslash{} \textbackslash{}text\{m\}\textasciicircum{}3\$ of helium gas at 1.00 atm and a temperature of 15.0\textbackslash{}degree C rises rapidly from ground level to an altitude at which the atmospheric pressure is only 0.900 atm (\textbackslash{}textbf\{figure\}). Assume the helium behaves like an ideal gas and the balloon’s ascent is too rapid to permit much heat exchange with the surrounding air.

\#\# Question

What is the change in internal energy of the helium as the balloon rises to the higher altitude?

\#\# Answer choices

- A: A: -1.65 \textbackslash{}times 10\textasciicircum{}\{9\}J

- B: B: -1.25 \textbackslash{}times 10\textasciicircum{}\{7\}J

- C: C: -1.95 \textbackslash{}times 10\textasciicircum{}\{7\}J

- D: D: -1.45 \textbackslash{}times 10\textasciicircum{}\{3\}J

\#\# Figure caption (auxiliary; use the image as primary evidence)

The image features two hot air balloons in flight. 

- The upper balloon is larger and in a warm orange color with a checker pattern. It has a visible basket beneath it.
- The lower balloon is slightly smaller, in a beige color with a stitched grid pattern. Its basket is also visible.
- Both balloons are floating above a simple landscape that includes hills, a few small buildings, a silo, and some trees, all depicted in a grayish tone. 

There is no text present in the image.

\#\# Dataset metadata

Ideal Gases and Kinetic Theory; Physical Model Grounding Reasoning, Numerical Reasoning

\paragraph{Recorded answer/context.}
B: B: -1.25 \textbackslash{}times 10\textasciicircum{}\{7\}J

\paragraph{Figure/image path.}
/root/proposal\_for\_physic/science-mango/phyx\_mini\_run/phyx\_data/test\_image/345.png

\paragraph{Formalization target.}
create a compiling Lean file with sorry bodies at `PhyXMiniProblems/problem\_phyx\_mini\_0345.lean`. The Lean declarations must preserve the physical quantities, dimensions or dimensional roles, figure labels, governing-law hypotheses, and final relation expressed by this problem.
Use Mathlib/Physlib names found through LeanExplore where available. If a physics API is missing, introduce faithful local abstractions rather than scalar placeholder aliases.

\begin{theorem}[Physics formalization target]
\label{thm:physics:phyx_mini_0345:target}
\lean{PhyXMiniProblems.ProblemPhyXMini0345.problem_phyx_mini_0345}
\uses{def:physics:phyx-mini-0345:phyxminiproblems-problemphyxmini0345-heliumballoonascent, def:physics:phyx-mini-0345:phyxminiproblems-problemphyxmini0345-matchesballoonscenario, def:physics:phyx-mini-0345:phyxminiproblems-problemphyxmini0345-matchesprimaryfigure, def:physics:phyx-mini-0345:phyxminiproblems-problemphyxmini0345-matchesproblemreadouts, def:physics:phyx-mini-0345:phyxminiproblems-problemphyxmini0345-hasphysicalballoonparameters, def:physics:phyx-mini-0345:phyxminiproblems-problemphyxmini0345-satisfiesrapidadiabaticidealgaslaws, def:physics:phyx-mini-0345:phyxminiproblems-problemphyxmini0345-internalenergychangeinjoules, def:physics:phyx-mini-0345:phyxminiproblems-problemphyxmini0345-isuniqueclosestdisplayedanswer}
The assigned autoformalize agent should translate this physics problem into Lean declarations in the covered file, with theorem and lemma proof bodies written as `by sorry`.
\end{theorem}
\begin{proof}
This is an autoformalization task, not a proof task. Produce statements that can later be proved by the physics prover without weakening the physical model.
\end{proof}

% --- Archon declaration topology begin ---
\section{Declaration topology}
\begin{definition}[Gas Volume]
  \label{def:physics:phyx-mini-0345:phyxminiproblems-problemphyxmini0345-gasvolume}
  \lean{PhyXMiniProblems.ProblemPhyXMini0345.GasVolume}
  A physical gas volume, independent of the selected system of units
\end{definition}
\begin{proof}
  This is the declared model definition of Gas Volume.
\end{proof}
\begin{definition}[volume In Cubic Meters]
  \label{def:physics:phyx-mini-0345:phyxminiproblems-problemphyxmini0345-volumeincubicmeters}
  \lean{PhyXMiniProblems.ProblemPhyXMini0345.volumeInCubicMeters}
  \uses{def:physics:phyx-mini-0345:phyxminiproblems-problemphyxmini0345-gasvolume}
  Read a gas volume in coherent SI cubic metres
\end{definition}
\begin{proof}
  This is the declared model definition of volume In Cubic Meters.
\end{proof}
\begin{definition}[pressure In Pascals]
  \label{def:physics:phyx-mini-0345:phyxminiproblems-problemphyxmini0345-pressureinpascals}
  \lean{PhyXMiniProblems.ProblemPhyXMini0345.pressureInPascals}
  Read a pressure in coherent SI pascals
\end{definition}
\begin{proof}
  This is the declared model definition of pressure In Pascals.
\end{proof}
\begin{definition}[energy In Joules]
  \label{def:physics:phyx-mini-0345:phyxminiproblems-problemphyxmini0345-energyinjoules}
  \lean{PhyXMiniProblems.ProblemPhyXMini0345.energyInJoules}
  Read an energy in coherent SI joules
\end{definition}
\begin{proof}
  This is the declared model definition of energy In Joules.
\end{proof}
\begin{definition}[temperature In Kelvins]
  \label{def:physics:phyx-mini-0345:phyxminiproblems-problemphyxmini0345-temperatureinkelvins}
  \lean{PhyXMiniProblems.ProblemPhyXMini0345.temperatureInKelvins}
  Physlib's Temperature.toReal is the absolute-temperature projection. In this model its selected scale is kelvin, as required by the ideal-gas law
\end{definition}
\begin{proof}
  This is the declared model definition of temperature In Kelvins.
\end{proof}
\begin{definition}[Balloon State]
  \label{def:physics:phyx-mini-0345:phyxminiproblems-problemphyxmini0345-balloonstate}
  \lean{PhyXMiniProblems.ProblemPhyXMini0345.BalloonState}
  The two endpoint states of the ascent
\end{definition}
\begin{proof}
  This is the declared model definition of Balloon State.
\end{proof}
\begin{definition}[Altitude Label]
  \label{def:physics:phyx-mini-0345:phyxminiproblems-problemphyxmini0345-altitudelabel}
  \lean{PhyXMiniProblems.ProblemPhyXMini0345.AltitudeLabel}
  The altitude labels supplied by the problem text
\end{definition}
\begin{proof}
  This is the declared model definition of Altitude Label.
\end{proof}
\begin{definition}[Balloon Gas]
  \label{def:physics:phyx-mini-0345:phyxminiproblems-problemphyxmini0345-balloongas}
  \lean{PhyXMiniProblems.ProblemPhyXMini0345.BalloonGas}
  The gas species carried by the research balloon
\end{definition}
\begin{proof}
  This is the declared model definition of Balloon Gas.
\end{proof}
\begin{definition}[Gas Model]
  \label{def:physics:phyx-mini-0345:phyxminiproblems-problemphyxmini0345-gasmodel}
  \lean{PhyXMiniProblems.ProblemPhyXMini0345.GasModel}
  The thermodynamic model applied to the balloon gas
\end{definition}
\begin{proof}
  This is the declared model definition of Gas Model.
\end{proof}
\begin{definition}[Ascent Protocol]
  \label{def:physics:phyx-mini-0345:phyxminiproblems-problemphyxmini0345-ascentprotocol}
  \lean{PhyXMiniProblems.ProblemPhyXMini0345.AscentProtocol}
  Qualitative speed of the ascent described in the source
\end{definition}
\begin{proof}
  This is the declared model definition of Ascent Protocol.
\end{proof}
\begin{definition}[Heat Exchange Regime]
  \label{def:physics:phyx-mini-0345:phyxminiproblems-problemphyxmini0345-heatexchangeregime}
  \lean{PhyXMiniProblems.ProblemPhyXMini0345.HeatExchangeRegime}
  Heat-exchange regime during the ascent
\end{definition}
\begin{proof}
  This is the declared model definition of Heat Exchange Regime.
\end{proof}
\begin{definition}[Figure Balloon]
  \label{def:physics:phyx-mini-0345:phyxminiproblems-problemphyxmini0345-figureballoon}
  \lean{PhyXMiniProblems.ProblemPhyXMini0345.FigureBalloon}
  The two balloon drawings in the supplied raster image
\end{definition}
\begin{proof}
  This is the declared model definition of Figure Balloon.
\end{proof}
\begin{definition}[Figure Balloon Color]
  \label{def:physics:phyx-mini-0345:phyxminiproblems-problemphyxmini0345-figureballooncolor}
  \lean{PhyXMiniProblems.ProblemPhyXMini0345.FigureBalloonColor}
  Fill colors visible in the supplied image
\end{definition}
\begin{proof}
  This is the declared model definition of Figure Balloon Color.
\end{proof}
\begin{definition}[Figure Balloon Pattern]
  \label{def:physics:phyx-mini-0345:phyxminiproblems-problemphyxmini0345-figureballoonpattern}
  \lean{PhyXMiniProblems.ProblemPhyXMini0345.FigureBalloonPattern}
  Envelope patterns visible in the supplied image
\end{definition}
\begin{proof}
  This is the declared model definition of Figure Balloon Pattern.
\end{proof}
\begin{definition}[Landscape Feature]
  \label{def:physics:phyx-mini-0345:phyxminiproblems-problemphyxmini0345-landscapefeature}
  \lean{PhyXMiniProblems.ProblemPhyXMini0345.LandscapeFeature}
  Landscape elements drawn below the two balloons
\end{definition}
\begin{proof}
  This is the declared model definition of Landscape Feature.
\end{proof}
\begin{definition}[Research Balloon Figure]
  \label{def:physics:phyx-mini-0345:phyxminiproblems-problemphyxmini0345-researchballoonfigure}
  \lean{PhyXMiniProblems.ProblemPhyXMini0345.ResearchBalloonFigure}
  \uses{def:physics:phyx-mini-0345:phyxminiproblems-problemphyxmini0345-balloonstate, def:physics:phyx-mini-0345:phyxminiproblems-problemphyxmini0345-figureballoon, def:physics:phyx-mini-0345:phyxminiproblems-problemphyxmini0345-figureballooncolor, def:physics:phyx-mini-0345:phyxminiproblems-problemphyxmini0345-figureballoonpattern, def:physics:phyx-mini-0345:phyxminiproblems-problemphyxmini0345-landscapefeature}
  Primary-image information is kept separately from thermodynamic quantities. The coordinates and rendered sizes are schematic scalar readouts, not physical altitudes or gas volumes
\end{definition}
\begin{proof}
  This is the declared model definition of Research Balloon Figure.
\end{proof}
\begin{definition}[Helium Balloon Ascent]
  \label{def:physics:phyx-mini-0345:phyxminiproblems-problemphyxmini0345-heliumballoonascent}
  \lean{PhyXMiniProblems.ProblemPhyXMini0345.HeliumBalloonAscent}
  \uses{def:physics:phyx-mini-0345:phyxminiproblems-problemphyxmini0345-gasvolume, def:physics:phyx-mini-0345:phyxminiproblems-problemphyxmini0345-balloonstate, def:physics:phyx-mini-0345:phyxminiproblems-problemphyxmini0345-altitudelabel, def:physics:phyx-mini-0345:phyxminiproblems-problemphyxmini0345-balloongas, def:physics:phyx-mini-0345:phyxminiproblems-problemphyxmini0345-gasmodel, def:physics:phyx-mini-0345:phyxminiproblems-problemphyxmini0345-ascentprotocol, def:physics:phyx-mini-0345:phyxminiproblems-problemphyxmini0345-heatexchangeregime, def:physics:phyx-mini-0345:phyxminiproblems-problemphyxmini0345-researchballoonfigure}
  The independent physical quantities of the balloon and its ascent. heliumAmountInMoles and molarGasConstantJoulesPerMoleKelvin are named scalar readouts in explicit units; pressure, volume, absolute temperature, heat, work, and internal energy remain physical quantities. No final energy-change value is stored here
\end{definition}
\begin{proof}
  This is the declared model definition of Helium Balloon Ascent.
\end{proof}
\begin{definition}[Matches Balloon Scenario]
  \label{def:physics:phyx-mini-0345:phyxminiproblems-problemphyxmini0345-matchesballoonscenario}
  \lean{PhyXMiniProblems.ProblemPhyXMini0345.MatchesBalloonScenario}
  \uses{def:physics:phyx-mini-0345:phyxminiproblems-problemphyxmini0345-heliumballoonascent}
  Qualitative facts stated in the prose, without any requested answer
\end{definition}
\begin{proof}
  This is the declared model definition of Matches Balloon Scenario.
\end{proof}
\begin{definition}[Matches Primary Figure]
  \label{def:physics:phyx-mini-0345:phyxminiproblems-problemphyxmini0345-matchesprimaryfigure}
  \lean{PhyXMiniProblems.ProblemPhyXMini0345.MatchesPrimaryFigure}
  \uses{def:physics:phyx-mini-0345:phyxminiproblems-problemphyxmini0345-figureballoon, def:physics:phyx-mini-0345:phyxminiproblems-problemphyxmini0345-landscapefeature, def:physics:phyx-mini-0345:phyxminiproblems-problemphyxmini0345-heliumballoonascent}
  Transcription of the supplied image. The upper orange checker-pattern balloon is rendered above and larger than the lower beige stitched-grid balloon; both have baskets. The image contains the listed landscape objects and no text. Its rendered-size inequality is not identified with a numerical gas volume
\end{definition}
\begin{proof}
  This is the declared model definition of Matches Primary Figure.
\end{proof}
\begin{definition}[Matches Problem Readouts]
  \label{def:physics:phyx-mini-0345:phyxminiproblems-problemphyxmini0345-matchesproblemreadouts}
  \lean{PhyXMiniProblems.ProblemPhyXMini0345.MatchesProblemReadouts}
  \uses{def:physics:phyx-mini-0345:phyxminiproblems-problemphyxmini0345-volumeincubicmeters, def:physics:phyx-mini-0345:phyxminiproblems-problemphyxmini0345-temperatureinkelvins, def:physics:phyx-mini-0345:phyxminiproblems-problemphyxmini0345-heliumballoonascent}
  Numerical readouts stated in the problem. The initial helium pressure is one standard atmosphere, the initial volume is 2000 m\textasciicircum{}3, and 15 degrees Celsius is encoded using the exact Celsius-to-kelvin offset 273.15. The only final-state datum is the ambient pressure 0.900 atm; equality of the helium and ambient endpoint pressures is a governing law below
\end{definition}
\begin{proof}
  This is the declared model definition of Matches Problem Readouts.
\end{proof}
\begin{definition}[Has Physical Balloon Parameters]
  \label{def:physics:phyx-mini-0345:phyxminiproblems-problemphyxmini0345-hasphysicalballoonparameters}
  \lean{PhyXMiniProblems.ProblemPhyXMini0345.HasPhysicalBalloonParameters}
  \uses{def:physics:phyx-mini-0345:phyxminiproblems-problemphyxmini0345-volumeincubicmeters, def:physics:phyx-mini-0345:phyxminiproblems-problemphyxmini0345-pressureinpascals, def:physics:phyx-mini-0345:phyxminiproblems-problemphyxmini0345-energyinjoules, def:physics:phyx-mini-0345:phyxminiproblems-problemphyxmini0345-temperatureinkelvins, def:physics:phyx-mini-0345:phyxminiproblems-problemphyxmini0345-heliumballoonascent}
  Positivity and nondegeneracy conditions for the thermodynamic model
\end{definition}
\begin{proof}
  This is the declared model definition of Has Physical Balloon Parameters.
\end{proof}
\begin{definition}[Satisfies Rapid Adiabatic Ideal Gas Laws]
  \label{def:physics:phyx-mini-0345:phyxminiproblems-problemphyxmini0345-satisfiesrapidadiabaticidealgaslaws}
  \lean{PhyXMiniProblems.ProblemPhyXMini0345.SatisfiesRapidAdiabaticIdealGasLaws}
  \uses{def:physics:phyx-mini-0345:phyxminiproblems-problemphyxmini0345-volumeincubicmeters, def:physics:phyx-mini-0345:phyxminiproblems-problemphyxmini0345-pressureinpascals, def:physics:phyx-mini-0345:phyxminiproblems-problemphyxmini0345-energyinjoules, def:physics:phyx-mini-0345:phyxminiproblems-problemphyxmini0345-temperatureinkelvins, def:physics:phyx-mini-0345:phyxminiproblems-problemphyxmini0345-balloonstate, def:physics:phyx-mini-0345:phyxminiproblems-problemphyxmini0345-heliumballoonascent}
  Governing endpoint laws for the idealized ascent: the flexible balloon is in mechanical equilibrium with the atmosphere at each endpoint; P V = n R T and U = (3/2) n R T hold for monatomic ideal helium; monatomic helium has adiabatic index gamma = 5/3; the adiabatic endpoint relation is T₂/T₁ = (P₂/P₁)\textasciicircum{}((gamma-1)/gamma); negligible heat exchange is idealized as Q = 0, and the first law is Delta U = Q - W\_by\_gas. These are generic physical relations between independently stored quantities.
\end{definition}
\begin{proof}
  This is the declared model definition of Satisfies Rapid Adiabatic Ideal Gas Laws.
\end{proof}
\begin{definition}[internal Energy Change In Joules]
  \label{def:physics:phyx-mini-0345:phyxminiproblems-problemphyxmini0345-internalenergychangeinjoules}
  \lean{PhyXMiniProblems.ProblemPhyXMini0345.internalEnergyChangeInJoules}
  \uses{def:physics:phyx-mini-0345:phyxminiproblems-problemphyxmini0345-energyinjoules, def:physics:phyx-mini-0345:phyxminiproblems-problemphyxmini0345-heliumballoonascent}
  The requested final-minus-initial internal-energy change, in joules
\end{definition}
\begin{proof}
  This is the declared model definition of internal Energy Change In Joules.
\end{proof}
\begin{definition}[Answer Choice]
  \label{def:physics:phyx-mini-0345:phyxminiproblems-problemphyxmini0345-answerchoice}
  \lean{PhyXMiniProblems.ProblemPhyXMini0345.AnswerChoice}
  Labels of the four displayed numerical answers
\end{definition}
\begin{proof}
  This is the declared model definition of Answer Choice.
\end{proof}
\begin{definition}[displayed Energy Change In Joules]
  \label{def:physics:phyx-mini-0345:phyxminiproblems-problemphyxmini0345-displayedenergychangeinjoules}
  \lean{PhyXMiniProblems.ProblemPhyXMini0345.displayedEnergyChangeInJoules}
  \uses{def:physics:phyx-mini-0345:phyxminiproblems-problemphyxmini0345-answerchoice}
  Internal-energy change in joules printed beside each answer label
\end{definition}
\begin{proof}
  This is the declared model definition of displayed Energy Change In Joules.
\end{proof}
\begin{definition}[recorded Dataset Answer]
  \label{def:physics:phyx-mini-0345:phyxminiproblems-problemphyxmini0345-recordeddatasetanswer}
  \lean{PhyXMiniProblems.ProblemPhyXMini0345.recordedDatasetAnswer}
  \uses{def:physics:phyx-mini-0345:phyxminiproblems-problemphyxmini0345-answerchoice}
  Dataset answer metadata, kept separate from every theorem premise
\end{definition}
\begin{proof}
  This is the declared model definition of recorded Dataset Answer.
\end{proof}
\begin{definition}[Is Unique Closest Displayed Answer]
  \label{def:physics:phyx-mini-0345:phyxminiproblems-problemphyxmini0345-isuniqueclosestdisplayedanswer}
  \lean{PhyXMiniProblems.ProblemPhyXMini0345.IsUniqueClosestDisplayedAnswer}
  \uses{def:physics:phyx-mini-0345:phyxminiproblems-problemphyxmini0345-answerchoice, def:physics:phyx-mini-0345:phyxminiproblems-problemphyxmini0345-displayedenergychangeinjoules}
  A choice is uniquely closest to an exact modeled energy change
\end{definition}
\begin{proof}
  This is the declared model definition of Is Unique Closest Displayed Answer.
\end{proof}
\begin{lemma}[final temperature from adiabatic expansion]
  \label{lem:physics:phyx-mini-0345:phyxminiproblems-problemphyxmini0345-final-temperature-from-adiabatic-expansion}
  \lean{PhyXMiniProblems.ProblemPhyXMini0345.final_temperature_from_adiabatic_expansion}
  \uses{def:physics:phyx-mini-0345:phyxminiproblems-problemphyxmini0345-temperatureinkelvins, def:physics:phyx-mini-0345:phyxminiproblems-problemphyxmini0345-heliumballoonascent, def:physics:phyx-mini-0345:phyxminiproblems-problemphyxmini0345-matchesballoonscenario, def:physics:phyx-mini-0345:phyxminiproblems-problemphyxmini0345-matchesproblemreadouts, def:physics:phyx-mini-0345:phyxminiproblems-problemphyxmini0345-hasphysicalballoonparameters, def:physics:phyx-mini-0345:phyxminiproblems-problemphyxmini0345-satisfiesrapidadiabaticidealgaslaws}
  The pressure ratio and monatomic adiabatic exponent determine T₂ = 288.15 * (0.9)\textasciicircum{}(2/5) K. This is an intermediate physical conclusion, not a premise
\end{lemma}
\begin{proof}
  The conclusion follows from the governing relations and figure readouts named in the statement of final temperature from adiabatic expansion.
\end{proof}
\begin{lemma}[internal energy change exact]
  \label{lem:physics:phyx-mini-0345:phyxminiproblems-problemphyxmini0345-internal-energy-change-exact}
  \lean{PhyXMiniProblems.ProblemPhyXMini0345.internal_energy_change_exact}
  \uses{def:physics:phyx-mini-0345:phyxminiproblems-problemphyxmini0345-heliumballoonascent, def:physics:phyx-mini-0345:phyxminiproblems-problemphyxmini0345-matchesballoonscenario, def:physics:phyx-mini-0345:phyxminiproblems-problemphyxmini0345-matchesproblemreadouts, def:physics:phyx-mini-0345:phyxminiproblems-problemphyxmini0345-hasphysicalballoonparameters, def:physics:phyx-mini-0345:phyxminiproblems-problemphyxmini0345-satisfiesrapidadiabaticidealgaslaws, def:physics:phyx-mini-0345:phyxminiproblems-problemphyxmini0345-internalenergychangeinjoules}
  Eliminating n R T₁ with the initial ideal-gas law gives the exact modeled energy change (3/2) * 101325 * 2000 * ((0.9)\textasciicircum{}(2/5) - 1) joules
\end{lemma}
\begin{proof}
  The conclusion follows from the governing relations and figure readouts named in the statement of internal energy change exact.
\end{proof}
% --- Archon declaration topology end ---
% --- Archon physics formalization source end ---
